% ===========================================================================
% OWNER: M3 (Zarpash Nasim); M2 adds runner test rows by PR
% GRADED BY: Form 3 #8 Testing plan & tool/risk disclosure (15)
% DRAW FROM: docs/05-Testing/evaluation-plan.md, fixture-suite.md, .github/workflows/ci.yml, eval/tests, runner/tests, evidence/tests, eval/reports/*.json, docs/evidence-m1/2/3.md
% NOTE: Result tables come from tables/generated/ (scripts/json_to_tex.py), never typed by hand
% Rules: mid-eval-report/CLAUDE.md. Do not edit a file you do not own.
% ===========================================================================
\chapter{Testing and Early Validation}

\section{Testing Strategy}
The evaluation plan (\texttt{docs/05-Testing/evaluation-plan.md}) treats testing, evaluation and
acceptance as three different questions, and this chapter only answers the first:
\textbf{testing} asks whether the implementation meets its own requirements (unit, integration
and system checks); \textbf{evaluation} asks whether the approach is better under stated budgets
(baseline comparisons, ablations); \textbf{acceptance} asks whether the intended users would
actually use it (named-participant tasks). Evaluation is partly covered in Chapter 4's
Acceptance Criteria (Table \ref{ch4:tab:acceptance}) via claim B11's baseline work; full
acceptance testing is genuinely not started (Section \ref{ch8:sec:acceptance}).

Strategy, concretely: every component is unit-tested in isolation with real, scriptable fakes
confined to \texttt{*/tests/} (never imported by production code --- CLAUDE.md \S4); the fixture
project (\texttt{fixtures/od-fixture}) then exercises the whole pipeline end to end against
seven cases with pre-registered, before-any-run ground truth
(\texttt{fixtures/od-fixture/ground\_truth.json}); and every one of those checks runs in CI on
every push and pull request, not only locally.

\section{Test Levels and Techniques}\label{ch8:sec:levels}
Table \ref{ch8:tab:ci} lists the five real CI jobs (\texttt{.github/workflows/ci.yml}), each
testing a different level, all green on the current \texttt{main} tip.

\begin{table}[H]\centering
\caption{CI jobs and what they test}\label{ch8:tab:ci}
\begin{tabular}{p{2.6cm}p{4.0cm}p{6.5cm}}
\toprule
\textbf{Job} & \textbf{Level} & \textbf{What it checks} \\
\midrule
\texttt{python-eval} & Unit (M3) & Every \texttt{eval/tests/test\_*.py} (85 tests), plus the
  yield report actually runs against the real benchmark manifest \\
\texttt{fixture-build} & Integration (fixture) & The fixture compiles; victims pass alone;
  victims fail after their real polluter; the N1 negative control fails alone --- catches a
  fixture change that would silently invalidate ground truth \\
\texttt{runner} & Unit + real JVM (M2) & Every \texttt{runner/tests/test\_*.py} (144 tests), with
  \texttt{FLAKETRACE\_REQUIRE\_JVM=1} so a missing JDK/Maven fails the job instead of silently
  skipping real-JVM tests \\
\texttt{evidence} & Unit + real \texttt{javap} (M1) & Every \texttt{evidence/tests/test\_*.py}
  (39 tests), against real compiled bytecode, same JVM-required discipline \\
\texttt{build-report} & Build & This document itself compiles (\texttt{latexmk -pdf main.tex}) \\
\bottomrule
\end{tabular}
\end{table}

\noindent All 268 unit tests (85 + 144 + 39) and the fixture/report builds pass on the current
\texttt{main} tip (CI run \texttt{38146685651}, all 5 jobs green --- see Appendix
\ref{appc:sec:ci}), confirmed by re-running the three suites locally while drafting this
chapter, not taken from a possibly-stale prior log. A project-wide discipline backs every one of them:
\texttt{eval/tests/fake\_runner.py}'s \texttt{FakeOrderRunner} and similar scriptable fakes live
\emph{only} under \texttt{*/tests/} and are never imported by production code (checked by
\texttt{grep}, not just asserted --- see \texttt{docs/evidence-m3.md}'s Phase 4 entry).
\section{Test Cases}\label{ch8:sec:cases}
Table \ref{ch8:tab:cases} lists all seven real cases in \texttt{fixtures/od-fixture}, each with
ground truth pre-registered in \texttt{ground\_truth.json} \textbf{before} any run against it
(\texttt{authored\_before\_any\_run: true} in the file itself) --- so a case's expected outcome
was never adjusted after seeing what the pipeline actually produced.

\begin{longtable}{p{0.8cm}p{2.1cm}p{8.0cm}p{2.3cm}}
\caption{The seven real fixture cases}\label{ch8:tab:cases}\\
\toprule
\textbf{Case} & \textbf{Category} & \textbf{What it tests} & \textbf{Expected} \\
\midrule
\endfirsthead
\toprule
\textbf{Case} & \textbf{Category} & \textbf{What it tests} & \textbf{Expected} \\
\midrule
\endhead
F1 & Single static-field polluter & A direct \texttt{getstatic}/\texttt{putstatic} write-read
  pair --- the baseline case & \texttt{VERIFIED} \\
F2 & System-property polluter & The victim reads the property \emph{inside a one-level helper
  call} (\texttt{FeatureFlags.isTurboEnabled()}), not its own test method --- tests that depth-2
  resolution finds evidence a depth-1 walk would miss & \texttt{VERIFIED} \\
F3 & Two required polluters & The victim fails only when \emph{both} predecessors have run;
  removing either alone makes it pass again --- tests \texttt{ddmin} minimisation, since a
  one-by-one search alone cannot find either polluter in isolation & \texttt{VERIFIED} \\
F4 & Order needs reshuffling & The polluter's class name sorts alphabetically \emph{after} the
  victim's, so the module's natural discovered order can never reproduce the bug --- tests
  ADR-008's shuffled-order search, the only path that can find this case at all &
  \texttt{VERIFIED} (via a shuffle) \\
N1 & Fails alone, deterministic & Always fails with no polluter run first --- must never be
  attributed to order dependence & \texttt{UNRESOLVED} (\texttt{VICTIM\_FAILS\_ALONE}) \\
N2 & Fails alone, intermittent & Fails \emph{sometimes} even when run alone, independent of any
  shared resource --- tests that the isolation check catches a recurring signature before
  reproduction counting is even consulted, distinguishing this from \texttt{NOT\_REPRODUCED} &
  \texttt{UNRESOLVED} (\texttt{VICTIM\_FAILS\_ALONE}) \\
N3 & Externally caused & Models a real bug whose actual cause (an environment variable
  FlakeTrace does not control) is outside anything the tool observes --- tests that the pipeline
  honestly reports it never reproduced the failure, rather than inventing a polluter or claiming
  the test is reliable & \texttt{UNRESOLVED} (\texttt{NOT\_REPRODUCED}) \\
\bottomrule
\end{longtable}

\textbf{A disclosed gap between the plan and what was built.}
\texttt{docs/05-Testing/fixture-suite.md} still describes an earlier, more ambitious ten-case
plan (F1--F10: reflection pollution, filesystem-path pollution, a nondeterministic victim, a
misleading failure signature, budget exhaustion, a non-OD control) that is \textbf{not} what
exists today. The suite that actually got built evolved differently during implementation ---
simpler on static/system-property coverage, but it gained F4 and N3 specifically to exercise
ADR-008's order-search and honest-non-reproduction paths, which the original ten-case plan never
anticipated. This is stated here rather than left as a silent mismatch between a planning
document and the real repository: the seven real cases are what Chapter 7's results and the
acceptance criteria in Chapter 4 are measured against, not the ten-case plan.

\section{Test Results}
% Generated from eval/reports/*.json — run: python3 mid-eval-report/scripts/json_to_tex.py
\IfFileExists{tables/generated/fixture_results.tex}{% GENERATED by mid-eval-report/scripts/json_to_tex.py from eval/reports/*.json - do not edit by hand.
{\footnotesize\setlength{\tabcolsep}{3pt}
\begin{longtable}{>{\raggedright\arraybackslash}p{0.04\textwidth}>{\raggedright\arraybackslash}p{0.18\textwidth}>{\raggedright\arraybackslash}p{0.16\textwidth}p{0.05\textwidth}p{0.1\textwidth}p{0.05\textwidth}p{0.04\textwidth}p{0.06\textwidth}>{\raggedright\arraybackslash}p{0.16\textwidth}}
\caption{Real diagnosis outcomes on the seeded fixture suite (20 replays per check, 95\% Wilson intervals)}\label{ch8:tab:fixture-results}\\
\toprule
Case & Victim & Outcome & Repro. & Wilson 95\% & Alone fail & Poll. & Order & Shared resource\\
\midrule
\endfirsthead
\toprule
Case & Victim & Outcome & Repro. & Wilson 95\% & Alone fail & Poll. & Order & Shared resource\\
\midrule
\endhead
F1 & Config\allowbreak{}Victim\allowbreak{}Test & VERIFIED & 20/20 & 0.84--1.00 & 0/20 & 1 & 3 $\to$ 2 & static-field: Config.mode \\
F2 & Feature\allowbreak{}Victim\allowbreak{}Test & VERIFIED & 20/20 & 0.84--1.00 & 0/20 & 1 & 6 $\to$ 2 & system-property: odfixture.turbo \\
F3 & Toggle\allowbreak{}Victim\allowbreak{}Test & VERIFIED & 20/20 & 0.84--1.00 & 0/20 & 2 & 15 $\to$ 3 & static-field: Toggles.flag\allowbreak{}A \\
F4 & Always\allowbreak{}Early\allowbreak{}Victim\allowbreak{}Test & VERIFIED & 20/20 & 0.84--1.00 & 0/20 & 1 & 3 $\to$ 2 & static-field: Late\allowbreak{}Flag.is\allowbreak{}Set \\
N1 & Negative\allowbreak{}Alone\allowbreak{}Fail\allowbreak{}Test & UNRESOLVED (VICTIM\_\allowbreak{}FAILS\_\allowbreak{}ALONE) & 20/20 & 0.84--1.00 & 20/20 & 0 & 9 $\to$ 1 & -- \\
N2 & Negative\allowbreak{}Flaky\allowbreak{}Test & UNRESOLVED (VICTIM\_\allowbreak{}FAILS\_\allowbreak{}ALONE) & 9/20 & 0.26--0.66 & 9/20 & 0 & 10 $\to$ 1 & -- \\
N3 & Env\allowbreak{}Dependent\allowbreak{}Negative\allowbreak{}Test & UNRESOLVED (NOT\_\allowbreak{}REPRODUCED) & 0/20 & 0.00--0.16 & 0/20 & 0 & 4 $\to$ 4 & -- \\
\bottomrule
\end{longtable}
}
}{\todo{run scripts/json\_to\_tex.py}}
\section{Acceptance Testing / Validation}\label{ch8:sec:acceptance}
\section{Edge and Failure Cases}
\section{Defects and Corrective Actions}
\section{Current Validation Status}
